\begin{center}
    {\Large\bfseries Устный тургор}
\end{center}

% Источник задачи: Турнир городов, международный (архив Торонто), весна 2012, старшая группа, сложный вариант (A-level), задача № 1; https://www.math.toronto.edu/oz/turgor/archives/TT2012S_SAsolutions.pdf
\problem{1}
В отряде охранников каждому назначено своё натуральное число. Для любых двух охранников отношение большего из назначенных им чисел к меньшему не меньше $3$. Охранник, которому назначено число $n$, дежурит $n$ дней подряд, затем $n$ дней подряд отдыхает, затем снова дежурит $n$ дней подряд и так далее. Начинать дежурства охранники могут в разные дни. Возможно ли устроить графики так, чтобы каждый день дежурил хотя бы один охранник?

\medskip

% Источник задачи: Турнир городов, международный (архив Торонто), весна 2012, старшая группа, сложный вариант (A-level), задача № 2; https://www.math.toronto.edu/oz/turgor/archives/TT2012S_SAsolutions.pdf
\problem{2}
Внутри круга отмечены $100$ точек, никакие три из которых не лежат на одной прямой. Докажите, что точки можно разбить на $50$ пар так, чтобы $50$ прямых, проходящих через точки каждой пары, попарно пересекались внутри круга.

\medskip

% Источник задачи: 62-я Международная математическая олимпиада, международная (Санкт-Петербург, Россия), 2021, шорт-лист, задача A3; https://www.imo-official.org/problems/IMO2021SL.pdf
\problem{3}
Пусть $n$ --- натуральное число. Найдите наименьшее возможное значение суммы
\[
  \left\lfloor\frac{a_1}{1}\right\rfloor+
  \left\lfloor\frac{a_2}{2}\right\rfloor+\dots+
  \left\lfloor\frac{a_n}{n}\right\rfloor,
\]
где $(a_1,a_2,\dots,a_n)$ пробегает все перестановки чисел $1,2,\dots,n$.

\medskip

% Источник задачи: 62-я Международная математическая олимпиада, международная (Санкт-Петербург, Россия), 2021, шорт-лист, задача G1; https://www.imo-official.org/problems/IMO2021SL.pdf
\problem{4}
Дан параллелограмм $ABCD$, причём $AC=BC$. Точка $P$ лежит на продолжении отрезка $AB$ за точку $B$. Описанная окружность треугольника $ACD$ вторично пересекает отрезок $PD$ в точке $Q$, а описанная окружность треугольника $APQ$ вторично пересекает отрезок $PC$ в точке $R$. Докажите, что прямые $CD$, $AQ$ и $BR$ проходят через одну точку.

\medskip

% Источник задачи: 62-я Международная математическая олимпиада, международная (Санкт-Петербург, Россия), 2021, шорт-лист, задача N3; https://www.imo-official.org/problems/IMO2021SL.pdf
\problem{5}
Найдите все натуральные числа $n$, обладающие следующим свойством: все $k$ натуральных делителей числа $n$ можно расположить в некотором порядке $d_1,d_2,\dots,d_k$ так, что при каждом $i=1,2,\dots,k$ сумма $d_1+d_2+\dots+d_i$ является точным квадратом.

\medskip

% Источник задачи: Канадская математическая олимпиада, Канада, 2020, задача № 4; https://w7.cms.math.ca/Competitions/CMO/archive/sol2020.pdf
\problem{6}
Пусть $S=\{1,4,8,9,16,\dots\}$ --- множество всех точных степеней натуральных чисел, то есть чисел вида $u^v$, где $u,v$ --- натуральные числа и $v\geq2$. Запишем элементы $S$ по возрастанию:
\[
  a_1<a_2<a_3<\dots.
\]
Докажите, что существует бесконечно много индексов $m$, для которых разность $a_{m+1}-a_m$ делится на $9999$.
